% 2026-10-23：按学习复现、探索尝试、优化验证、迭代调试、整板尝试重组。
\documentclass[stage2,fontset=fandol]{ydzjreport}
\usepackage{subcaption,longtable}

\stagetwosetup
{第 5 个逻辑月 · 整板尝试}
{冯旭东、陈思文、张昊洋}
{复合微通道分区冷却}
{【待确认】逻辑月 5 对应的自然月}
{整板布局；网络连接；供液；集成尝试；验证准备}
{承接小块探索与适用边界研究，尝试整板布局、单元连接及供液组织，区分已有几何准备与尚待求解的整板性能。}
\pagestyle{empty}   % 去掉页眉与页脚（含页码）
\RemoveFromHook{enddocument}[ydzjreport] % 去掉末页的作者与格式规范收尾说明


\makeatletter
\newcommand{\submissionnumber}[1]{\def\@currentlabel{#1}}
\makeatother
\begin{document}

\makestagetwohead

% unified-five-part-outline 2026-10-23
\section{本月计划}

\subsection{任务目标}

本月将局部小块研究推进到整板尝试，先组织布局、连接与供液，再检查是否具备流热计算和控制验证条件。整板几何准备不等于整板性能已验证；月四分界线未补齐时，已有布局只作为候选方案。

\subsection{前置条件}


\par\noindent \begingroup\submissionnumber{1}\phantomsection\label{tab:pre}\endgroup 表~1 请参见附件表~1。\par




\section{研究过程}

\subsection{整板布局与连接尝试}

\subsubsection{已有几何准备及其边界}
中期资料记录了 $40\,\mathrm{mm}\times40\,\mathrm{mm}$ 整板、$10\times10$ 单元及40 mm条带装配；其中低、中、高负荷分别为70、27、3格。此处沿用已有几何与热源映射，不把候选配置写成小块分界线已经验证的产物。各版本的主体、节距和连接尺寸需对应实际文件核查。

当前整板尚缺完整的流体域提取、共享拓扑、网格及求解证据。尝试顺序是核对几何和热源，再明确拓扑与供回液边界，随后建立流体域与网格，检查守恒与收敛，最后进行性能对照。

\subsubsection{连接与供液的基本约束}

\paragraph{单元到整板的条件变化}
单元级的仿真结果是在
"独立、充分供液"的假设下得到的；一旦连成整板：

\begin{itemize}
  \item \textbf{并联}时，高阻力支路可能\textbf{供液不足}，
        原本温度最低的流道反而得不到足够流量；
  \item \textbf{串联}时，上游吸热会\textbf{抬高下游入口水温}，
        使下游小块的冷却能力下降。
\end{itemize}

因此已有单元级结果\textbf{不能直接搬到整板}。

\paragraph{连接段阻力与偏流}
小块之间的过渡与分流通道会引入额外压降，
且容易在局部产生回流与偏流。连接段设计的评价，应同时看
\textbf{入口速度分布的均匀性}与\textbf{连接段压降}。

% ==================================================================

\subsubsection{连接段设计目标与评价方法}

已确定的\textbf{设计目标}：优化连接段截面及分流布局，
减少局部回流、流动突变和偏流，使各小块入口的速度分布尽量均匀、
流动整体保持稳定。

已确定的\textbf{评价方法}：结合入口速度分布、支路流量及连接段压降评价效果，
并按各区冷却需求合理分配流量。

\subsection{后续验证方案的设计}

\subsubsection{验证一：分区方案是否有效}

\textbf{要回答的问题}：采用差异化结构的分区冷板，
是否真的改善了温度分布，而不是仅仅改变了平均值。

\textbf{指标}：
\begin{itemize}
  \item 整板\textbf{最高温度}——是否低于允许上限；
  \item 整板\textbf{温度均匀性}——用同一指定表面的
        最高温与最低温之差，或温度标准差表征；
  \item 各区\textbf{实际温度}与设计预期的偏差。
\end{itemize}

\textbf{判据}：与均匀布置方案对比，在相同泵功下比较温度均匀性，改善幅度及合格阈值待明确。

\subsubsection{验证二：流动分配是否合理}

\textbf{要回答的问题}：并联或混合网络下，各小块是否都得到了设计流量。

\textbf{指标}：
\begin{itemize}
  \item 各\textbf{支路流量}相对设计值的偏差；
  \item \textbf{连接段压降}及其占整板压降的比例；
  \item 各小块\textbf{入口速度分布}的均匀性。
\end{itemize}

\textbf{判据}：无支路出现明显供液不足；入口速度分布无严重偏流。

\subsubsection{验证三：能效是否可接受}

\textbf{指标}：整板泵功、综合性能指标、峰值热阻。

\textbf{注意}：综合指标必须与单元级使用\textbf{同一基准}，
否则整板与单元不可比。

\subsubsection{验证四：是否优于单流道}

\textbf{要回答的问题}：分区复合方案相对单一流道，收益是否成立。

\textbf{指标}：相同热载荷下的最高温度、温度均匀性、泵功三项对比。

\textbf{注意}：前述分析说明\textbf{单元排序不能直接搬到整板}，
因此本项对比必须用整板数据，不能引用单元级结论。

\subsubsection{验证五：控制是否可行}

\textbf{要回答的问题}：控制方向的两个待辨识函数能否落地。

\textbf{要辨识的对象}：
\begin{itemize}
  \item 执行器到流速的关系（泵转速、阀开度 $\rightarrow$ 流速）；
  \item 流速与水温到芯片温度的关系。
\end{itemize}

\textbf{所需条件}：泵的型号与流量范围、是否已加装水阀、
有无流量计、温度采样频率，
以及\textbf{各小块之间是否独立供液}（决定控制架构）。

【缺数据：上述控制实验条件】
【待确认8】控制方向所需数据的提供方式与格式。

% ==================================================================

\section{完成情况}

\subsection{任务完成状态}


\par\noindent \begingroup\submissionnumber{2}\phantomsection\label{tab:tasks}\endgroup 表~2 请参见附件表~2。\par




\subsection{可以确定的}

已具备开展整板尝试的局部资料、部分几何与连接方向；整板性能应以独立求解和同条件对照判断。

\subsection{尚不能确定的}

现有布局是否最优、整板温度均匀性、实际泵送代价和多行协同控制收益均尚缺系统级数据。

\section{问题与初步对策}

\subsection{遇到的主要问题}

局部排序无法直接用于整板；热源覆盖面积与总功率需要一致；连接阻力、支路分流和下游入口升温需要整板或条带计算；控制执行器尚缺实测辨识。

\subsection{下一步对策}

先以候选几何跑通整体流程；记录连接压降与支路流量，作为结构和控制共用输入；整板采用独立网格检验，并对物性适用温区、质量和能量平衡进行核查。

\section{下个月计划}


\par\noindent \begingroup\submissionnumber{3}\phantomsection\label{tab:plan}\endgroup 表~3 请参见附件表~3。\par




\begin{thebibliography}{9}
\bibitem{proposal} 大创项目申报书：研究内容与实施方案. 2026.
\bibitem{data} 四种流道数据与折线图. 项目计算资料, 2026.
\bibitem{udf} 项目组. water\_properties\_refprop\_fit.c：液态水物性拟合与 Fluent UDF 源码. 项目计算资料, 2026.
\end{thebibliography}

\end{document}
